\documentclass[11pt,a4paper]{article}
\usepackage{brand}

% Provider colours, used consistently in every figure.
\colorlet{cclef}{amber}
\colorlet{cjev}{blue}
\colorlet{claya}{teal}
\colorlet{cclm}{navy}

% Scientific notation for p-values: \sci{1.1}{-5}
\newcommand\sci[2]{$#1\times10^{#2}$}

\brandsetup{
  title    = {Four Decision Models, One Protocol},
  subtitle = {A paired evaluation of CLEF-Flash, Jev, Laya, and CLM on intent triage, reranking, knowledge, routing, and tool selection},
  category = {Evaluation},
  version  = {v1.0},
  doctype  = {paper},
}

\begin{document}
\maketitle

\begin{abstractbox}
Decision models answer typed questions, such as ``which of these options applies?'', with a probability
for every option instead of free text. That makes them natural routers, triage classifiers, and tool
selectors inside agentic systems. This report compares four of them, CLEF-Flash, Jev, Laya, and CLM,
by sending each one identical requests through the same SystemOne decision API and the same sequential
client. Five suites cover Banking77 intent triage (154 messages, 77 options), Amazon ESCI reranking
(50 queries, 10 products each), MMLU (200 questions from ten subjects), a curated 90-prompt
model-routing set, and the full 200-case BFCL v3 multiple tool-selection split. All 2,776 measured
requests completed without errors or retries. CLEF-Flash and Jev form a clear top tier. CLEF-Flash
reaches 93.5\% on Banking77 against Jev's 80.5\% (exact McNemar test, Holm-adjusted
$p = 1.1\times10^{-5}$), while Jev leads on the other four suites by margins that are not significant
(Holm $p$ between 0.25 and 1.0). Both top models significantly outperform Laya and CLM on every suite.
Laya is the fastest model (8--20 ms median latency) and reaches 94.0\% on tool selection, but trails
far behind on MMLU and routing. The zero-shot CLM checkpoint collapses onto two intents on Banking77
and lands within 0.008 nDCG of a random ordering on ESCI. Alongside the headline numbers, the report
gives Wilson intervals, paired bootstrap tests, calibration metrics, and an error analysis showing that
two of the five routing errors made by the top models coincide with an alternative reference label.
\end{abstractbox}

\tableofcontents
\newpage

% =====================================================================
\section{Introduction}\label{sec:intro}
% =====================================================================

Agentic systems make many small decisions before they generate any text: which queue a support
ticket belongs to, which tool a request needs, and how much model capacity a prompt deserves. A
decision model answers those questions directly. It receives a \emph{state} and a set of named
\emph{questions}, each listing its options, and returns a choice together with a probability for
every option. The four models studied here, CLEF-Flash \cite{cloudflare2026clefflash},
Jev \cite{typesafe2026jev}, Laya \cite{convai2026laya}, and CLM \cite{contrastivelm2026clm}, all
expose this interface through the same SystemOne endpoint. A shared protocol makes a like-for-like
comparison possible at the level of requests rather than prompts.

Comparisons between such models are easy to get wrong. Option order, question wording, client
timeouts, warm server caches, and GPU sharing can each shift accuracy by several points, and a single
unpaired accuracy figure hides how often two models disagree on the same case. This study holds every
one of those factors fixed and measures what remains, using paired tests that compare the models case
by case.

The report makes four contributions:
\begin{enumerate}
  \item A controlled comparison in which every provider receives the same state, questions, option
        text, and option order from one shared client, with the three local models served one at a
        time on the same GPU (\cref{sec:method}).
  \item Five suites totalling 694 measured requests per provider, with every per-case output and
        run manifest saved, so each number in this report can be recomputed without calling a model
        (\cref{sec:experiments}).
  \item Paired significance tests with family-wise correction, interval estimates, calibration
        metrics, and explicit chance baselines, including a random-permutation floor for ESCI and a
        uniform-random baseline for BFCL (\cref{sec:results}).
  \item An error analysis that separates genuine routing failures from label ambiguity in the curated
        routing set (\cref{sec:routing}).
\end{enumerate}

\Cref{sec:related} positions the work, \cref{sec:method} defines the protocol, task formulations,
and statistics, \cref{sec:experiments} describes the data and serving setup, \cref{sec:results}
reports the findings, and \cref{sec:discussion} discusses limitations before \cref{sec:conclusion}
concludes. The next section starts with the work this evaluation builds on.

% =====================================================================
\section{Related Work}\label{sec:related}
% =====================================================================

Three lines of work shape this evaluation: model routing, the public benchmarks reused here, and the
statistics of comparing classifiers. This section covers each in turn.

\paragraph{Routing and cascades.} FrugalGPT \cite{chen2023frugalgpt} cascades calls across language
models of increasing cost and stops once an answer is judged reliable. RouteLLM \cite{ong2024routellm}
learns routers between a strong and a weak model from preference data. The routing suite in this report
asks a narrower question: given an explicit rubric, can a decision model name the smallest capability
tier that can complete a request?

\paragraph{Benchmarks.} Banking77 \cite{casanueva2020banking77} is a fine-grained intent detection set
with 77 banking intents. The Shopping Queries Dataset \cite{reddy2022esci} labels query and product
pairs as Exact, Substitute, Complement, or Irrelevant (ESCI). MMLU \cite{hendrycks2021mmlu} measures
knowledge across 57 subjects with four-way multiple choice. The Berkeley Function Calling Leaderboard
\cite{gorilla2024bfcl}, from the Gorilla project \cite{patil2023gorilla}, evaluates tool calls; its
\emph{multiple} split offers several candidate functions per request. Every suite here is a reduced or
reframed version of its source (\cref{sec:suites}), so the scores are diagnostic and not comparable to
leaderboard numbers.

\paragraph{Evaluation methodology.} Dietterich \cite{dietterich1998tests} recommends McNemar's test
\cite{mcnemar1947} for comparing two classifiers on the same test set, and Holm's procedure
\cite{holm1979} controls the family-wise error rate across many pairwise tests. Wilson intervals
\cite{wilson1927} remain accurate near 0\% and 100\%, where the normal approximation fails. Ranking
quality uses nDCG \cite{jarvelin2002ndcg} with the exponential gain popularised by RankNet
\cite{burges2005ranknet}, compared through a paired bootstrap \cite{efron1993bootstrap}. Calibration
uses log loss, the Brier score \cite{brier1950}, and expected calibration error
\cite{naeini2015ece,guo2017calibration}. Finally, language models are known to be sensitive to the
order in which options appear \cite{zheng2024selectors,pezeshkpour2024order}, which motivates the
deterministic option shuffle described next.

With that context in place, the next section defines exactly what each model was asked and how the
answers were scored.

% =====================================================================
\section{Methodology}\label{sec:method}
% =====================================================================

The methodology has four parts: the request protocol shared by all models, the way each benchmark is
cast into that protocol, the controls that keep the comparison fair, and the statistics used to
compare models. \Cref{fig:harness} shows how the parts connect.

\begin{figure}[htbp]
  \begin{brandflow}
    \node[brand node] (data) at (0,0) {Prepared datasets\\[1pt]{\scriptsize\mdseries 5 suites, seed 20261009}};
    \node[brand blue node] (build) at (4.7,0) {Request builder\\[1pt]{\scriptsize\mdseries same state, options, order}};
    \node[brand dark node] (client) at (9.4,0) {Shared client\\[1pt]{\scriptsize\mdseries 120 s timeout, 3 attempts}};
    \node[brand amber node] (clef) at (0,-2.3) {CLEF-Flash\\[1pt]{\scriptsize\mdseries local, GPU 1}};
    \node[brand amber node] (jev) at (3.13,-2.3) {Jev\\[1pt]{\scriptsize\mdseries hosted API}};
    \node[brand amber node] (laya) at (6.27,-2.3) {Laya\\[1pt]{\scriptsize\mdseries local, GPU 1}};
    \node[brand amber node] (clm) at (9.4,-2.3) {CLM\\[1pt]{\scriptsize\mdseries local, GPU 1}};
    \begin{scope}[on background layer]
      \node[brand group, fit=(clef)(jev)(laya)(clm)] (prov) {};
    \end{scope}
    \node[brand node] (rows) at (2.35,-4.7) {Per-case rows\\[1pt]{\scriptsize\mdseries 2,776 results, run manifests}};
    \node[brand green node] (stats) at (7.05,-4.7) {Paired analysis\\[1pt]{\scriptsize\mdseries McNemar, Holm, bootstrap}};
    \draw[brand arrow] (data) -- (build);
    \draw[brand arrow] (build) -- (client);
    \draw[brand arrow] (client) -- (client |- prov.north);
    \draw[brand arrow] (prov.south -| rows) -- (rows);
    \draw[brand arrow] (rows) -- (stats);
  \end{brandflow}
  \caption{Evaluation harness. One request builder and one client serve all four providers; the three
  local models run one at a time on the same GPU, and every answer is stored as a per-case row before
  any statistic is computed.}\label{fig:harness}
\end{figure}

\subsection{The SystemOne decision protocol}\label{sec:protocol}

Every provider accepts the same request shape at \code{POST /v1/systemone}. The \code{state} field
holds the input as a JSON object. The \code{questions} field maps question names to question
definitions; this study uses only the \code{choice} type, whose \code{criteria} map each option name to
a short description. \Cref{code:request} shows an MMLU request in this form.

\begin{brandcode}[lang=json, caption={An MMLU request in SystemOne form. Options are shown in A--D order; the harness shuffles them per case before sending\label{code:request}}]
{
  "state": {
    "subject": "world religions",
    "question": "What is the most important festival for Sikhs?"
  },
  "questions": {
    "answer": {
      "type": "choice",
      "instructions": "Select the one correct answer to the multiple-choice question.",
      "criteria": {"A": "Holi", "B": "Divali", "C": "Baisakhi Day", "D": "Hola Mahalla"}
    }
  },
  "model": "clef-flash"
}
\end{brandcode}

The response carries one answer per question, with the chosen option and a probability for each
option, plus the serving model's identifier and token usage. \Cref{code:response} reproduces a real
response from the routing run.

\begin{brandcode}[lang=json, caption={A saved Jev response to routing case small-001, a title-case conversion request\label{code:response}}]
{
  "model": "jev-1.13.0",
  "answers": {
    "route": {
      "type": "choice", "choice": "small", "confidence": 1.0,
      "probabilities": {"small": 1.0, "medium": 0.0, "large": 0.0}
    }
  },
  "usage": {"input_tokens": 524, "output_tokens": 38}
}
\end{brandcode}

Because the probabilities are part of the protocol, every suite can be scored for both accuracy and
calibration. The next subsection explains how each benchmark is mapped onto this request shape.

\subsection{Task formulation}\label{sec:suites}

Each suite becomes one SystemOne request per case, summarised in \cref{tab:suites}. The formulations
are deliberately plain: option names and descriptions come straight from the source data, so no
provider benefits from prompt engineering.

\begin{brandtable}[The five evaluation suites\label{tab:suites}]{L L C C L}{Suite, Source, Cases, Options, Primary metric}
  Banking77 triage   & \cite{casanueva2020banking77} & 154 & 77 & Accuracy \\
  ESCI reranking     & \cite{reddy2022esci}          & 50 queries & 4 per product & nDCG@10 \\
  MMLU               & \cite{hendrycks2021mmlu}      & 200 & 4 & Accuracy \\
  Model routing      & curated                       & 90 & 3 & Accuracy \\
  BFCL v3 multiple   & \cite{gorilla2024bfcl}        & 200 & 2--4 & Accuracy \\
\end{brandtable}

\paragraph{Banking77 triage.} The state is the customer message. A single \code{queue} question
offers all 77 intents, each described by its label with underscores replaced by spaces (for example,
``card arrival''). Secondary metric: macro-F1.

\paragraph{MMLU.} The state holds the subject and the question; the four answer texts are the options
A--D.

\paragraph{ESCI reranking.} One request per query asks a separate four-way relevance question for each
of its ten products. Each product $j$ receives an expected-relevance score from the returned
probabilities $p_j$:
\begin{equation}\label{eq:score}
  s_j = \sum_{g \in \mathcal{G}} r(g)\, p_j(g),
\end{equation}
where $\mathcal{G}$ is the set of ESCI grades and $r$ maps Exact, Substitute, Complement, and Irrelevant
to 3, 2, 1, and 0. Products are ranked by descending $s_j$, and the ranking is scored with
\begin{equation}\label{eq:ndcg}
  \text{nDCG@10} = \frac{\text{DCG}}{\text{IDCG}}, \qquad
  \text{DCG} = \sum_{i=1}^{10} \frac{2^{r_{(i)}} - 1}{\log_2(i + 1)},
\end{equation}
where $r_{(i)}$ is the gold grade of the product at rank $i$ and IDCG is the DCG of the ideal order.
The exponential gain gives Exact, Substitute, Complement, and Irrelevant gains of 7, 3, 1, and 0.
The secondary metric is the mean reciprocal rank (MRR) of the first Exact product.

\paragraph{Model routing.} The state is the user prompt. A single \code{route} question asks for the
\emph{minimum} capability tier (\code{small}, \code{medium}, or \code{large}) that can reliably complete
the whole request, with four tie-breaking rules: prefer the lowest sufficient tier; high consequence
raises a request to \code{large}; several otherwise-medium deliverables together are \code{large}; and
input length alone does not raise the tier when the work is mechanical.

\paragraph{BFCL tool selection.} The state is the flattened conversation. The options are the exact
candidate function names, each described by its documentation string and its JSON parameter schema.
Only the selected function name is scored, not its arguments.

These formulations fix what is asked. The next subsection fixes how it is asked, so that the only
remaining difference between runs is the model.

\subsection{Fairness controls}\label{sec:fairness}

Five controls keep the comparison fair.

\begin{itemize}
  \item \textbf{Identical requests.} Every provider receives the same state, questions, option text,
        and option order. The only provider-specific fields are each vendor's model identifier and,
        for Laya, its documented token-budget fields (\cref{sec:providers}).
  \item \textbf{Deterministic option order.} For Banking77, MMLU, ESCI, and BFCL the option order is
        shuffled per case, with a seed derived from the case identifier $c$:
        \begin{equation}\label{eq:shuffle}
          \pi_c = \operatorname{Shuffle}\bigl(\mathcal{O}_c;\ \operatorname{seed} = \operatorname{SHA-256}(\texttt{20261009:}\,c)\bigr).
        \end{equation}
        The order is therefore random with respect to the gold label but identical across providers.
        The routing question keeps a fixed order, discussed in \cref{sec:threats}.
  \item \textbf{One client.} All requests go through one sequential HTTP client with a 120 s timeout
        and up to three attempts with exponential backoff (0.5 s, then 1 s). Latency is measured on
        the client for the successful attempt.
  \item \textbf{Cold-cache protection.} Each provider run starts with one warm-up request that is not
        part of any dataset, so no server-side cache can answer a measured case.
  \item \textbf{One GPU, one model at a time.} CLEF-Flash, Laya, and CLM ran sequentially on the same
        NVIDIA RTX 4090, each from a freshly restarted server. Jev ran remotely in the same session.
\end{itemize}

Each provider run followed the same procedure:

\begin{steps}
  \stepitem{Start the server fresh}{Restart the provider's container on the benchmark GPU so no state carries over.}
  \stepitem{Record the environment}{Append the endpoint, model identifier, extra request fields, retry policy, health response, and GPU memory to the run manifest.}
  \stepitem{Warm up}{Send one request that belongs to no dataset.}
  \stepitem{Run the suites}{Send every pending case sequentially and append each answer as a row the moment it arrives.}
  \stepitem{Stop the server}{Shut the container down before the next local model starts.}
\end{steps}

These controls make the per-case outputs directly comparable. The final methodological question is
how to decide whether an observed difference is real.

\subsection{Statistical analysis}\label{sec:stats}

All comparisons are paired: the same cases go to every model, so the tests compare models case by case
rather than through independent samples.

\paragraph{Interval estimates.} Accuracy on Banking77 and MMLU is reported with a 95\% Wilson interval,
\begin{equation}\label{eq:wilson}
  \frac{\hat{p} + \frac{z^2}{2n} \pm z\sqrt{\frac{\hat{p}(1-\hat{p})}{n} + \frac{z^2}{4n^2}}}{1 + \frac{z^2}{n}},
  \qquad z = 1.96,
\end{equation}
and tested against chance with a one-sided exact binomial test (chance is 1/77 for Banking77 and 1/4 for
MMLU).

\paragraph{Paired accuracy tests.} For two models on the same $n$ cases, let $b$ count cases only the
first model gets right and $c$ cases only the second gets right. The exact two-sided McNemar test is
\begin{equation}\label{eq:mcnemar}
  p = \min\left\{1,\; 2 \sum_{k=0}^{\min(b,c)} \binom{b+c}{k}\, 2^{-(b+c)}\right\}.
\end{equation}

\paragraph{Paired ranking tests.} For ESCI, the per-query nDCG difference between two models is
resampled over queries $B = 20{,}000$ times with replacement. The 2.5th and 97.5th percentiles of the
resampled mean differences $\Delta^{*}_{b}$ give the confidence interval, and the two-sided $p$-value is
\begin{equation}\label{eq:bootp}
  p_{\text{boot}} = \min\left\{1,\; 2\min\left(\frac{1}{B}\sum_{b=1}^{B} \mathbf{1}[\Delta^{*}_{b} \le 0],\;
  \frac{1}{B}\sum_{b=1}^{B} \mathbf{1}[\Delta^{*}_{b} \ge 0]\right)\right\}.
\end{equation}

\paragraph{Multiple comparisons.} Each suite has $m = 6$ provider pairs. With the raw $p$-values sorted
as $p_{(1)} \le \dots \le p_{(m)}$, Holm's adjusted values are
\begin{equation}\label{eq:holm}
  \tilde{p}_{(i)} = \max_{j \le i}\, \min\bigl\{1,\; (m - j + 1)\, p_{(j)}\bigr\}.
\end{equation}
A difference is called significant when $\tilde{p} < 0.05$.

\paragraph{Calibration.} For $N$ cases with gold labels $y_i$, options $\mathcal{K}_i$, and returned
probabilities $p_i$,
\begin{align}
  \text{LogLoss} &= -\frac{1}{N}\sum_{i=1}^{N} \log \max\bigl(p_i(y_i),\, 10^{-12}\bigr), \label{eq:logloss}\\
  \text{Brier}   &= \frac{1}{N}\sum_{i=1}^{N}\sum_{k \in \mathcal{K}_i} \bigl(p_i(k) - \mathbf{1}[k = y_i]\bigr)^2, \label{eq:brier}\\
  \text{ECE}     &= \sum_{b=1}^{10} \frac{|B_b|}{N}\,\bigl|\operatorname{acc}(B_b) - \operatorname{conf}(B_b)\bigr|, \label{eq:ece}
\end{align}
where ECE groups cases into ten equal-width bins $B_b$ by their top-class probability. ECE is computed
for routing and tool selection; log loss and Brier are computed for every single-choice suite. Every
provider returned a complete probability distribution on every case, so no case drops out of these
metrics.

With the protocol and statistics defined, the next section describes the data and the serving setup
that produced the results.

% =====================================================================
\section{Experimental Setup}\label{sec:experiments}
% =====================================================================

This section covers the prepared datasets, the four serving configurations, and the execution record
of the run.

\subsection{Datasets and sampling}\label{sec:data}

The three public subsets were sampled once with seed 20261009 and saved, together with SHA-256 hashes
of the source files (\cref{tab:sources}). BFCL uses the complete official split. The routing set was
written for this study.

\begin{brandtable}[Data sources and sampling. Hashes are the first eight hex digits of each source file's SHA-256\label{tab:sources}]{L L L C}{Suite, Source files, Sampling, SHA-256}
  Banking77 & \code{mteb/banking77} test \cite{hf2026banking77} & 2 messages per intent, 77 intents & 9575f636 \\
  ESCI & \code{tasksource/esci} test, shard 0 \cite{hf2026esci} & 50 of 1,855 eligible US queries; 10 products each, at least one Exact and one Irrelevant & 01b0da5f \\
  MMLU & \code{cais/mmlu} all, test \cite{hf2026mmlu} & 20 questions from each of 10 fixed subjects & 74a41822 \\
  BFCL & v3 multiple inputs and answers \cite{gorilla2024bfcl} & all 200 cases & aef16815, 244e00ce \\
  Routing & curated & 30 prompts per tier, 87 task categories & n/a \\
\end{brandtable}

Two properties of the samples matter for interpretation. First, 218 of the 500 sampled ESCI products
are labelled Exact, which raises the score of a random ordering (\cref{sec:reranking}). Second, the
routing file carries an alternative reference label for every prompt, which disagrees with the gold
label on 5 of the 90 prompts; \cref{sec:routing} uses it to examine routing errors. BFCL cases offer
two candidates in 79 cases, three in 85, and four in 36, so choosing uniformly at random would score
$(79/2 + 85/3 + 36/4)/200 = 38.4\%$.

\subsection{Providers and serving configuration}\label{sec:providers}

Each local model ran through its vendor's own serving code with the configuration its documentation
recommends (\cref{tab:providers}).

\begin{brandtable}[Serving configuration. The model column shows the identifier each server reported in its responses\label{tab:providers}]{l L l l}{Provider, Serving stack, Model, Extra fields}
  CLEF-Flash & Qwen3.5-9B backbone with a joint schema head, loaded by the model's own release code; transformers 5.10.2, PyTorch 2.11 & \code{clef-flash} & none \\
  Jev & TypeSafe hosted API & \code{jev-1.13.0} & none \\
  Laya & \code{laya.serve} 0.3.21 with its Router choosing among English, multilingual, and typed-decision checkpoints & \code{laya-rl-agent} & token budget \\
  CLM & \code{clm-serve} 0.1.0 heads over Qwen3-8B \cite{qwen2025qwen3} last-token embeddings served by vLLM 0.19.0 \cite{kwon2023vllm} in bfloat16 & \code{clm-latest} & none \\
\end{brandtable}

Laya's request omits a model identifier, so its Router selects the checkpoint; it chose the English
checkpoint for all of these English inputs. Laya also received \code{max_len} = 2048 and
\code{head_max_len} = 1024. At Laya's suggested \code{head_max_len} of 512, Banking77's 77 options
collapse to 72 distinct spans; 1024 is the smallest budget tested at which Laya reports no merged
options. CLM ran its projection heads on the CPU with its action cache set to zero, behind a vLLM
encoder limited to 2,048 tokens.

\subsection{Execution record}\label{sec:execution}

The run manifests record the whole evaluation on 11 October 2026 between 20:24 and 20:33 UTC, in the
order Jev, CLM, CLEF-Flash, Laya. Each provider answered 404 public-suite cases, 90 routing prompts,
and 200 BFCL cases: 694 measured requests per provider and 2,776 in total. Every request succeeded on
its first attempt. The GPU snapshots taken at the start of each run show 18,675--19,647 MiB in use on
the benchmark GPU for CLEF-Flash, 21,579--21,679 MiB for CLM, and 3,375--4,359 MiB for Laya. A second
RTX 4090 in the same host held 20,636 MiB for an unrelated workload throughout.

\begin{successbox}[Complete, error-free run]
All 2,776 measured requests returned a valid answer on the first attempt, and every answer included a
probability for every option. No case is missing from any metric in this report.
\end{successbox}

With a complete and clean set of per-case answers, the next section turns to what they show.

% =====================================================================
\section{Results}\label{sec:results}
% =====================================================================

The results lead with the headline comparison and its significance, then examine each suite, and
close with calibration and latency.

\noindent\metric{Banking77, CLEF-Flash}{93.5\%}[+13.0 vs Jev]\hfill\metric{BFCL, Jev}{100\%}[200/200]\hfill\metric{Median latency, Laya}{8.4 ms}

\begin{insightbox}
Two tiers emerge. CLEF-Flash and Jev are statistically indistinguishable on four of five suites; the
one significant gap between them is fine-grained intent triage, where CLEF-Flash leads by 13 points.
Both are significantly better than Laya and CLM on every suite.
\end{insightbox}

\subsection{Headline comparison}\label{sec:headline}

\Cref{tab:headline} lists the primary metric for every suite, and \cref{fig:primary} plots it against
each suite's chance or random baseline.

\begin{brandtable}[Primary metric by suite. Accuracies in percent; best value in bold\label{tab:headline}]{L c c c c c}{Suite (metric), CLEF-Flash, Jev, Laya, CLM, Chance}
  Banking77 (accuracy)    & \textbf{93.5} & 80.5 & 45.5 & 1.3 & 1.3 \\
  ESCI (nDCG@10)          & 0.9139 & \textbf{0.9183} & 0.8553 & 0.8012 & 0.7940 \\
  MMLU (accuracy)         & 92.5 & \textbf{93.5} & 34.5 & 29.5 & 25.0 \\
  Routing (accuracy)      & 96.7 & \textbf{97.8} & 47.8 & 35.6 & 33.3 \\
  BFCL (accuracy)         & 98.5 & \textbf{100.0} & 94.0 & 65.5 & 38.4 \\
\end{brandtable}

\begin{figure}[htbp]
  \centering
  \begin{tikzpicture}
    \begin{axis}[brand plot, ybar=1.5pt, bar width=7pt, width=\linewidth, height=6.6cm,
        ymin=0, ymax=1.08, enlarge x limits=0.12,
        symbolic x coords={Banking77,ESCI,MMLU,Routing,BFCL}, xtick=data,
        ylabel={Primary metric}, area legend,
        legend style={at={(0.5,-0.13)}, anchor=north, legend columns=-1,
          /tikz/every even column/.append style={column sep=8pt}}]
      \addplot[fill=cclef, draw=cclef] coordinates {(Banking77,0.9351) (ESCI,0.9139) (MMLU,0.925) (Routing,0.9667) (BFCL,0.985)};
      \addplot[fill=cjev, draw=cjev] coordinates {(Banking77,0.8052) (ESCI,0.9183) (MMLU,0.935) (Routing,0.9778) (BFCL,1.0)};
      \addplot[fill=claya, draw=claya] coordinates {(Banking77,0.4545) (ESCI,0.8553) (MMLU,0.345) (Routing,0.4778) (BFCL,0.94)};
      \addplot[fill=cclm, draw=cclm] coordinates {(Banking77,0.0130) (ESCI,0.8012) (MMLU,0.295) (Routing,0.3556) (BFCL,0.655)};
      \legend{CLEF-Flash, Jev, Laya, CLM}
      \draw[densely dashed, brandgray, line width=1pt] ([xshift=-19pt]axis cs:Banking77,0.013) -- ([xshift=19pt]axis cs:Banking77,0.013);
      \draw[densely dashed, brandgray, line width=1pt] ([xshift=-19pt]axis cs:ESCI,0.794) -- ([xshift=19pt]axis cs:ESCI,0.794);
      \draw[densely dashed, brandgray, line width=1pt] ([xshift=-19pt]axis cs:MMLU,0.25) -- ([xshift=19pt]axis cs:MMLU,0.25);
      \draw[densely dashed, brandgray, line width=1pt] ([xshift=-19pt]axis cs:Routing,0.3333) -- ([xshift=19pt]axis cs:Routing,0.3333);
      \draw[densely dashed, brandgray, line width=1pt] ([xshift=-19pt]axis cs:BFCL,0.384) -- ([xshift=19pt]axis cs:BFCL,0.384);
    \end{axis}
  \end{tikzpicture}
  \caption{Primary metric by suite. Dashed lines mark chance accuracy (Banking77, MMLU, routing), the
  uniform-random baseline for BFCL, and the random-permutation nDCG@10 for ESCI.}\label{fig:primary}
\end{figure}

The picture is consistent across suites: the top two models sit near the ceiling everywhere except
Banking77, Laya sits between the tiers, and CLM stays close to its baseline on three suites. Whether
those visual gaps are real is the subject of the next subsection.

\subsection{Pairwise significance}\label{sec:pairwise}

Each suite has six provider pairs, tested with \cref{eq:mcnemar} (or \cref{eq:bootp} for ESCI) and
Holm-adjusted within the suite (\cref{eq:holm}). \Cref{tab:clefjev} isolates the top-tier comparison.

\begin{brandtable}[CLEF-Flash vs Jev. ``Only'' columns count cases one model answers correctly and the other does not; the effect is CLEF-Flash minus Jev\label{tab:clefjev}]{l c c L c c}{Suite, CLEF only, Jev only, Effect, $p$, $p_{\text{Holm}}$}
  Banking77 & 21 & 1  & +13.0 pts & \sci{1.1}{-5} & \sci{1.1}{-5} \\
  ESCI      & -- & -- & $-0.0044$ nDCG [$-0.042$, $0.028$] & 0.840 & 0.840 \\
  MMLU      & 9  & 11 & $-1.0$ pts & 0.824 & 0.824 \\
  Routing   & 2  & 3  & $-1.1$ pts & 1.000 & 1.000 \\
  BFCL      & 0  & 3  & $-1.5$ pts & 0.250 & 0.250 \\
\end{brandtable}

Banking77 is decisive: of the 22 cases on which the two models disagree in correctness, CLEF-Flash wins
21. Everywhere else the discordant counts are small and balanced, and the ESCI interval comfortably
spans zero. A larger sample could still reveal a difference on BFCL, where Jev wins all three
discordant cases, but three cases cannot establish one.

\begin{brandtable}[Laya vs CLM. The effect is Laya minus CLM\label{tab:layaclm}]{l c c L c c}{Suite, Laya only, CLM only, Effect, $p$, $p_{\text{Holm}}$}
  Banking77 & 70 & 2  & +44.2 pts & \sci{1.1}{-18} & \sci{3.3}{-18} \\
  ESCI      & -- & -- & $+0.0541$ nDCG [$0.008$, $0.104$] & 0.019 & 0.037 \\
  MMLU      & 46 & 36 & +5.0 pts & 0.320 & 0.641 \\
  Routing   & 25 & 14 & +12.2 pts & 0.108 & 0.216 \\
  BFCL      & 64 & 7  & +28.5 pts & \sci{1.3}{-12} & \sci{5.0}{-12} \\
\end{brandtable}

\Cref{tab:layaclm} shows that Laya is significantly better than CLM on Banking77, ESCI, and BFCL, while
MMLU and routing do not separate them. All 20 cross-tier comparisons (CLEF-Flash or Jev against Laya or
CLM, on each of the five suites) are significant after correction; the weakest is CLEF-Flash against
Laya on BFCL (11 against 2 discordant cases, Holm $p = 0.045$). \Cref{app:pairwise} lists all 24 McNemar
tests, and \cref{tab:boot} gives the six ESCI comparisons.

\begin{infobox}[Reading the pairwise tables]
McNemar's test uses only the discordant cases. Two models with similar accuracy can still differ
significantly if one wins most disagreements, and a large accuracy gap built on few disagreements can
fail to reach significance. The ``only'' columns show both counts so the reader can judge the evidence
directly.
\end{infobox}

The pairwise tests settle the ranking. The remaining subsections explain where in each suite the
models succeed and fail.

\subsection{Intent triage and knowledge}\label{sec:triage}

\Cref{tab:ci} adds 95\% Wilson intervals and chance tests to the Banking77 and MMLU accuracies.

\begin{brandtable}[Accuracy with 95\% Wilson intervals and one-sided binomial tests against chance\label{tab:ci}]{l C C C C}{Provider, Banking77 [95\% CI], $p$ vs 1/77, MMLU [95\% CI], $p$ vs 1/4}
  CLEF-Flash & 93.5 [88.5, 96.4] & \sci{3.0}{-257} & 92.5 [88.0, 95.4] & \sci{8.4}{-92} \\
  Jev        & 80.5 [73.6, 86.0] & \sci{6.2}{-203} & 93.5 [89.2, 96.2] & \sci{5.6}{-95} \\
  Laya       & 45.5 [37.8, 53.3] & \sci{2.3}{-88}  & 34.5 [28.3, 41.3] & 0.0017 \\
  CLM        & 1.3 [0.4, 4.6]    & 0.596           & 29.5 [23.6, 36.2] & 0.084 \\
\end{brandtable}

On Banking77, the CLEF-Flash and Jev intervals do not overlap, which agrees with the McNemar result.
Macro-F1 tells the same story (0.932 for CLEF-Flash, 0.791 for Jev, 0.438 for Laya, 0.001 for CLM).
CLM is the striking case: it predicted only two of the 77 intents, \code{exchange_via_app} for 108
messages and \code{card_acceptance} for the other 46, and its accuracy is indistinguishable from chance
($p = 0.596$).

On MMLU, the top two models are within one point of each other and both clear 88\% at the lower
interval bound. Laya is above chance ($p = 0.0017$) but more than 50 points behind; CLM does not reach
significance against chance ($p = 0.084$). The subject breakdown in \cref{app:mmlu} shows CLEF-Flash's
weakest subject is anatomy (75\%), while Jev's is business ethics (80\%); with 20 questions per subject,
these differences are within noise.

Both suites reward knowledge of the label space. The reranking suite asks a different question: how
well the probabilities order products.

\subsection{Reranking}\label{sec:reranking}

\Cref{tab:esci} reports the ranking metrics alongside per-product label accuracy, and \cref{tab:boot}
gives the paired bootstrap comparisons.

\begin{brandtable}[ESCI ranking and per-product label metrics. The random row is the mean over 20,000 random orderings\label{tab:esci}]{l C C C C}{Provider, nDCG@10, MRR (Exact), Label accuracy, Label macro-F1}
  CLEF-Flash      & 0.9139 & 0.8545 & \textbf{0.568} & \textbf{0.532} \\
  Jev             & \textbf{0.9183} & \textbf{0.8670} & 0.546 & 0.489 \\
  Laya            & 0.8553 & 0.7730 & 0.428 & 0.304 \\
  CLM             & 0.8012 & 0.6329 & 0.282 & 0.189 \\
  Random ordering & 0.7940 & 0.6277 & --    & --    \\
\end{brandtable}

\begin{brandtable}[Paired bootstrap on per-query nDCG@10 (20,000 resamples). A $p$ of $<10^{-4}$ means no resample crossed zero\label{tab:boot}]{L c c c c}{Pair (A minus B), Mean difference, 95\% CI, $p$, $p_{\text{Holm}}$}
  CLEF-Flash minus Jev  & $-0.0044$ & [$-0.0422$, $0.0282$] & 0.840  & 0.840 \\
  CLEF-Flash minus Laya & $+0.0586$ & [$0.0210$, $0.0965$]  & 0.0015 & 0.0045 \\
  CLEF-Flash minus CLM  & $+0.1127$ & [$0.0754$, $0.1520$]  & $<10^{-4}$ & $<10^{-4}$ \\
  Jev minus Laya        & $+0.0630$ & [$0.0278$, $0.0988$]  & 0.0006 & 0.0024 \\
  Jev minus CLM         & $+0.1171$ & [$0.0721$, $0.1664$]  & $<10^{-4}$ & $<10^{-4}$ \\
  Laya minus CLM        & $+0.0541$ & [$0.0082$, $0.1036$]  & 0.0185 & 0.037 \\
\end{brandtable}

The two top models rank products almost equally well, yet they label them differently: CLEF-Flash
assigns the exact ESCI grade more often (56.8\% against 54.6\%), while Jev's probabilities produce a
marginally better ordering. Label accuracy stays below 57\% for every model, which reflects how
ambiguous the Substitute and Complement grades are, even though the rankings are good.

\begin{warningbox}[Read nDCG against its floor]
A random ordering already scores nDCG@10 = 0.794 and MRR = 0.628 here, because 44\% of the sampled
products are Exact and many queries have several. The usable range is therefore roughly 0.79 to 1.0, not
0 to 1. CLM's 0.801 sits within 0.008 of random, and the 0.06 gap between Laya and the top tier is
larger than it looks.
\end{warningbox}

Reranking compares models on how they spread probability. Model routing instead asks them to judge the
difficulty of a request, which turns out to separate the tiers sharply.

\subsection{Model routing}\label{sec:routing}

\Cref{tab:routing} and \cref{fig:routing} break routing accuracy down by expected tier.

\begin{brandtable}[Routing accuracy, macro-F1, per-tier recall, and the distribution of predicted tiers (small/medium/large; 30 of each is ideal)\label{tab:routing}]{l C C C C C C}{Provider, Accuracy, Macro-F1, Small, Medium, Large, Predicted}
  CLEF-Flash & 96.7 & 0.967 & 1.000 & 0.900 & 1.000 & 30/27/33 \\
  Jev        & \textbf{97.8} & \textbf{0.978} & 1.000 & 0.933 & 1.000 & 32/28/30 \\
  Laya       & 47.8 & 0.438 & 0.567 & 0.767 & 0.100 & 23/60/7 \\
  CLM        & 35.6 & 0.249 & 0.900 & 0.167 & 0.000 & 68/21/1 \\
\end{brandtable}

\begin{figure}[htbp]
  \centering
  \begin{tikzpicture}
    \begin{axis}[brand plot, ybar=1.5pt, bar width=9pt, width=0.85\linewidth, height=6cm,
        ymin=0, ymax=1.08, enlarge x limits=0.25,
        symbolic x coords={Small,Medium,Large}, xtick=data,
        ylabel={Recall}, area legend,
        legend style={at={(0.5,-0.15)}, anchor=north, legend columns=-1,
          /tikz/every even column/.append style={column sep=8pt}}]
      \addplot[fill=cclef, draw=cclef] coordinates {(Small,1.0) (Medium,0.9) (Large,1.0)};
      \addplot[fill=cjev, draw=cjev] coordinates {(Small,1.0) (Medium,0.9333) (Large,1.0)};
      \addplot[fill=claya, draw=claya] coordinates {(Small,0.5667) (Medium,0.7667) (Large,0.1)};
      \addplot[fill=cclm, draw=cclm] coordinates {(Small,0.9) (Medium,0.1667) (Large,0.0)};
      \legend{CLEF-Flash, Jev, Laya, CLM}
    \end{axis}
  \end{tikzpicture}
  \caption{Routing recall by expected tier. Laya funnels most prompts to \emph{medium} and CLM to
  \emph{small}, so neither recognises large requests.}\label{fig:routing}
\end{figure}

The weaker models fail in opposite directions. Laya sends 60 of 90 prompts to \emph{medium} and
recognises only 3 of 30 large requests; CLM sends 68 prompts to \emph{small} and recognises none. Both
patterns would be costly in production: Laya would under-provision difficult work, and CLM would
under-provision almost everything.

The top models make only five errors between them, all on medium prompts (\cref{tab:routeerr}).
CLEF-Flash errs upward three times, routing medium work to \emph{large}; Jev errs downward twice,
routing it to \emph{small}.

\begin{brandtable}[Every routing error made by CLEF-Flash or Jev, with the alternative reference label\label{tab:routeerr}]{l L l l l l}{Case, Task, Gold, Alternative, Model, Predicted}
  medium-011 & Migration notice with impact and rollback contact & medium & medium & Jev & small \\
  medium-019 & Secure password-reset flow with token expiry & medium & large & CLEF-Flash & large \\
  medium-020 & Regex for ticket IDs, with explanation & medium & small & Jev & small \\
  medium-021 & Idempotent, resumable nightly CSV import & medium & medium & CLEF-Flash & large \\
  medium-028 & Minor framework upgrade with rollback plan & medium & medium & CLEF-Flash & large \\
\end{brandtable}

\begin{experimentbox}[Errors or label ambiguity?]
Two of the five errors, medium-019 and medium-020, are prompts where the alternative reference label
agrees with the model rather than the gold label. The password-reset prompt is security-sensitive, which
the rubric's high-consequence rule could reasonably push to \emph{large}; the regex prompt is close to a
single mechanical step. Switching to the alternative labels would not favour either model, though: they
rescue one error per model but contradict four answers each model currently gets right, so CLEF-Flash
would score 93.3\% and Jev 94.4\%. The gold labels are curated judgments, and differences of one or two
cases between the top models should be read in that light.
\end{experimentbox}

Routing judgements depend on understanding what a request entails. Tool selection depends mostly on
matching a request to a function description, and the gap between the tiers narrows accordingly.

\subsection{Tool selection}\label{sec:tools}

\Cref{tab:bfcl} reports BFCL accuracy and calibration.

\begin{brandtable}[BFCL v3 multiple tool selection (200 cases)\label{tab:bfcl}]{l C C C C C}{Provider, Accuracy, Macro-F1, Log loss, Brier, ECE}
  CLEF-Flash     & 98.5  & 0.984 & 0.0500 & 0.0129 & 0.0449 \\
  Jev            & \textbf{100.0} & \textbf{1.000} & \textbf{0.0027} & \textbf{0.0006} & \textbf{0.0025} \\
  Laya           & 94.0  & 0.936 & 0.2289 & 0.1132 & 0.1138 \\
  CLM            & 65.5  & 0.635 & 0.9995 & 0.5263 & 0.1578 \\
  Uniform random & 38.4  & --    & --     & --     & --     \\
\end{brandtable}

Jev selects the right function in all 200 cases with almost no probability on the alternatives.
CLEF-Flash's three errors are all near-misses between closely related functions:
\code{store_price} against \code{on_sale} in a video-games module, Biot--Savart against Amp\`ere's
law in an electromagnetism module, and a wildlife population-growth function against a broader
ecological impact analysis. Laya's 94.0\% is its best result in the study and the suite where it comes closest to
the top tier. CLM is far above the uniform baseline (65.5\% against 38.4\%), which shows its encoder
carries useful signal even where its zero-shot heads fail elsewhere.

Accuracy alone does not say whether a model's probabilities can be trusted as confidence scores. The
next subsection examines that directly.

\subsection{Calibration}\label{sec:calibration}

\Cref{tab:calib} collects log loss and Brier scores for the single-choice suites, and \cref{fig:calib}
compares each model's mean top-class probability with its accuracy on routing and BFCL.

\begin{brandtable}[Calibration by suite (lower is better). ECE is reported for routing and BFCL only\label{tab:calib}]{L c c c c}{Suite and metric, CLEF-Flash, Jev, Laya, CLM}
  Banking77 log loss & \textbf{0.274} & 1.998 & 8.766 & 5.880 \\
  Banking77 Brier    & \textbf{0.094} & 0.291 & 0.954 & 1.123 \\
  MMLU log loss      & \textbf{0.244} & 0.456 & 1.393 & 2.575 \\
  MMLU Brier         & \textbf{0.111} & 0.112 & 0.758 & 1.050 \\
  Routing log loss   & 0.226 & \textbf{0.067} & 0.909 & 1.707 \\
  Routing Brier      & 0.093 & \textbf{0.038} & 0.550 & 0.888 \\
  Routing ECE        & 0.140 & \textbf{0.044} & 0.149 & 0.358 \\
  BFCL log loss      & 0.050 & \textbf{0.003} & 0.229 & 0.999 \\
\end{brandtable}

\begin{figure}[htbp]
  \centering
  \begin{tikzpicture}
    \begin{axis}[brand plot, width=0.72\linewidth, height=7cm,
        xmin=0.4, xmax=1.02, ymin=0.3, ymax=1.03,
        xlabel={Mean top-class probability}, ylabel={Accuracy},
        legend pos=south east]
      \addplot[densely dashed, brandgray, domain=0.4:1.02, samples=2, forget plot] {x};
      \addplot[only marks, mark=*, mark size=3pt, color=cclef] coordinates {(0.826,0.967)};
      \addplot[only marks, mark=*, mark size=3pt, color=cjev] coordinates {(0.962,0.978)};
      \addplot[only marks, mark=*, mark size=3pt, color=claya] coordinates {(0.498,0.478)};
      \addplot[only marks, mark=*, mark size=3pt, color=cclm] coordinates {(0.714,0.356)};
      \addplot[only marks, mark=square*, mark size=2.6pt, color=cclef, forget plot] coordinates {(0.956,0.985)};
      \addplot[only marks, mark=square*, mark size=2.6pt, color=cjev, forget plot] coordinates {(0.998,1.0)};
      \addplot[only marks, mark=square*, mark size=2.6pt, color=claya, forget plot] coordinates {(0.833,0.94)};
      \addplot[only marks, mark=square*, mark size=2.6pt, color=cclm, forget plot] coordinates {(0.791,0.655)};
      \legend{CLEF-Flash, Jev, Laya, CLM}
    \end{axis}
  \end{tikzpicture}
  \caption{Mean top-class probability against accuracy, computed from the saved per-case rows. Circles
  are routing, squares are BFCL. Points above the dashed diagonal are underconfident; points below it
  are overconfident.}\label{fig:calib}
\end{figure}

Three patterns stand out. CLEF-Flash has the best log loss and Brier score on Banking77 and MMLU, but
it is \emph{underconfident} on routing: its mean confidence is 0.83 while its accuracy is 0.97, which
explains its routing ECE of 0.140. Jev shows the opposite failure on Banking77: its log loss of 2.0
comes almost entirely from its 30 errors, where the gold intent receives a mean negative log
probability of 9.8 (a probability of roughly $6\times10^{-5}$), against 0.12 on its correct answers.
CLM is \emph{overconfident} on both routing (0.71 mean confidence at 36\% accuracy) and BFCL (0.79 at
66\%). Laya's Banking77 log loss of 8.8 is consistent with its checkpoint's own warning that the
calibration temperature for questions with eleven or more options is invalid and was clamped.

Calibration describes how much to trust a single answer. Latency describes what each answer costs, and
the last subsection reports it.

\subsection{Latency}\label{sec:latency}

\Cref{tab:latency} gives client-side median and 95th-percentile latency for every suite.

\begin{brandtable}[Client-side latency in milliseconds, median / 95th percentile. Jev includes the network round trip to its hosted API\label{tab:latency}]{l C C C C}{Suite, CLEF-Flash, Jev, Laya, CLM}
  Banking77 & 322.6 / 323.7 & 327.2 / 411.1 & 16.3 / 17.1 & 25.7 / 26.6 \\
  ESCI      & 309.0 / 332.1 & 312.3 / 410.2 & 19.5 / 21.0 & 66.8 / 76.2 \\
  MMLU      & 69.1 / 87.5   & 315.5 / 410.6 & 8.4 / 9.6   & 61.3 / 68.9 \\
  Routing   & 86.2 / 86.6   & 309.4 / 400.2 & 9.0 / 9.4   & 28.2 / 30.2 \\
  BFCL      & 102.5 / 140.4 & 312.0 / 408.4 & 8.9 / 9.6   & 68.7 / 74.8 \\
\end{brandtable}

Laya is an order of magnitude faster than every other model on every suite. CLEF-Flash's latency tracks
request size: about 70--100 ms for the four-option and tool questions, but over 300 ms for Banking77's
77 options and ESCI's ten questions per request. Jev's latency is nearly constant at about 310 ms, which
suggests that network time dominates; it should not be read as a measure of Jev's compute cost.

\begin{conclusionbox}[Results in brief]
CLEF-Flash and Jev are tied on knowledge, ranking, routing, and tool selection, and CLEF-Flash is
clearly better at fine-grained intent triage. Laya is fast and strong at tool selection but weak at
judging difficulty. The zero-shot CLM checkpoint is near chance on three suites. Calibration and latency
add trade-offs that accuracy alone hides.
\end{conclusionbox}

These findings raise the question of why the tiers separate where they do, and how far the conclusions
generalise. The discussion takes up both.

% =====================================================================
\section{Discussion}\label{sec:discussion}
% =====================================================================

This section interprets the tier structure, sets out the threats to validity, turns the findings into
practical guidance, and lists the open questions.

\subsection{Where the tiers separate}\label{sec:interpretation}

The suites differ in how much they demand beyond matching text to an option. BFCL tool selection is
largely lexical: the request names an operation and one candidate's documentation describes it, so even
Laya reaches 94\%. Banking77 is also matching, but across 77 fine-grained and partly overlapping intents,
which separates CLEF-Flash from Jev. MMLU requires knowledge that is not present in the request, and
routing requires judging how difficult and consequential a task is; on these, Laya and CLM fall to
between 30\% and 48\% while the top two stay above 92\%.

\begin{insightbox}
Matching a request to a described option is close to solved for the top models, and partly within reach
of a small, fast model. Judging difficulty or recalling knowledge is where model capacity still decides
the outcome.
\end{insightbox}

\subsection{Threats to validity}\label{sec:threats}

Several factors limit how far these results generalise.

\begin{itemize}
  \item \textbf{Diagnostic subsets.} Banking77 uses two messages per intent, ESCI ten sampled products
        per query, MMLU ten of 57 subjects, and BFCL scores only the function name. None of these are
        leaderboard scores.
  \item \textbf{Single run.} Each case was answered once per model, so run-to-run variation is not
        measured. Jev is a hosted model whose version can change; its Banking77 accuracy was 79.2\% in
        an earlier two-model run against 80.5\% here, while CLEF-Flash reproduced exactly.
  \item \textbf{Fixed routing order.} The routing question always lists \emph{small}, \emph{medium},
        \emph{large} in that order. Every provider saw the same order, so the comparison is fair, but
        position bias \cite{zheng2024selectors,pezeshkpour2024order} is not controlled on that suite.
  \item \textbf{Curated routing labels.} The routing gold labels are single curated judgments, and the
        alternative labels disagree on 5 of 90 prompts (\cref{sec:routing}).
  \item \textbf{Vendor configurations.} Laya ran with a larger token budget than its default; default
        settings would score lower on Banking77. CLM was evaluated as the zero-shot checkpoint, while its
        model card attributes its published agentic results to fine-tuned heads. Reproducing the card's
        own example gave the same decision but a probability of 0.988 instead of the documented 0.939, so
        a small numerical difference in the encoder cannot be ruled out.
  \item \textbf{Recorded versions.} The public-suite rows predate per-row model identifiers, so Jev's
        version there (jev-1.13.0 in the routing and BFCL rows) is inferred from running in the same
        session. Model-weight revisions and two unpinned Python dependencies in the CLEF-Flash image
        were not recorded.
  \item \textbf{Latency.} Jev's figures include network time, and the local figures depend on this GPU,
        on request size, and on a host whose second GPU carried an unrelated workload.
\end{itemize}

\subsection{Practical guidance}\label{sec:guidance}

Within those limits, the results support a simple decision framework for choosing among these four
models.

\begin{decision}
  \decisionitem{the task has many fine-grained labels, as in intent triage}{CLEF-Flash}{it leads Jev by 13 points on Banking77, the only significant top-tier gap}
  \decisionitem{the decision needs knowledge or a difficulty judgement}{CLEF-Flash or Jev}{neither is significantly better on MMLU, ESCI, or routing}
  \decisionitem{tool selection must fit a tight latency budget}{Laya}{94.0\% at a 9 ms median, though 4.5 points below CLEF-Flash (Holm $p = 0.045$)}
  \decisionitem{probabilities will gate downstream actions}{check calibration per task}{CLEF-Flash is underconfident on routing and Jev overconfident on Banking77 errors}
  \decisionitem{only the zero-shot CLM checkpoint is available}{train task heads first}{it is at or near chance on Banking77, MMLU, and routing}
\end{decision}

\subsection{Future work}\label{sec:future}

Four extensions would sharpen these conclusions:
\begin{enumerate}
  \item Shuffle the routing options per case and rerun, to measure position bias directly.
  \item Repeat the hosted-model runs on several days to estimate version drift.
  \item Evaluate full Banking77 and MMLU test sets, which would narrow the intervals in \cref{tab:ci}.
  \item Score BFCL arguments as well as function names, and add a fine-tuned CLM for comparison.
\end{enumerate}

% =====================================================================
\section{Conclusion}\label{sec:conclusion}
% =====================================================================

Four decision models received identical requests through one protocol and one client, on five suites
and 2,776 measured requests, with every per-case answer saved. CLEF-Flash and Jev form a clear top tier
and are statistically tied on four suites; CLEF-Flash's 13-point lead on fine-grained intent triage is
the one significant difference between them. Laya trades accuracy for speed and is competitive only on
tool selection, and the zero-shot CLM checkpoint is near chance on three suites. Paired tests, interval
estimates, and explicit baselines turned a table of accuracies into claims that can be defended, and the
saved per-case rows let anyone recompute or extend every one of them.

% =====================================================================
\appendix
% =====================================================================

\section{All pairwise McNemar tests}\label{app:pairwise}

The two tables below list every McNemar test behind \cref{sec:pairwise}. ``Left'' and ``right'' follow
alphabetical provider order, as in the saved result files.

\begin{brandtable}[Exact McNemar tests for Banking77 and MMLU, Holm-adjusted within each suite\label{tab:mcnemar-public}]{l l C C C C}{Suite, {Left, right}, Left only, Right only, $p$, $p_{\text{Holm}}$}
  Banking77 & CLEF-Flash, CLM  & 142 & 0   & \sci{3.6}{-43} & \sci{2.2}{-42} \\
  Banking77 & CLEF-Flash, Jev  & 21  & 1   & \sci{1.1}{-5}  & \sci{1.1}{-5} \\
  Banking77 & CLEF-Flash, Laya & 75  & 1   & \sci{2.0}{-21} & \sci{8.2}{-21} \\
  Banking77 & CLM, Jev         & 0   & 122 & \sci{3.8}{-37} & \sci{1.9}{-36} \\
  Banking77 & CLM, Laya        & 2   & 70  & \sci{1.1}{-18} & \sci{3.3}{-18} \\
  Banking77 & Jev, Laya        & 59  & 5   & \sci{9.0}{-13} & \sci{1.8}{-12} \\
  MMLU      & CLEF-Flash, CLM  & 129 & 3   & \sci{1.4}{-34} & \sci{7.0}{-34} \\
  MMLU      & CLEF-Flash, Jev  & 9   & 11  & 0.824          & 0.824 \\
  MMLU      & CLEF-Flash, Laya & 119 & 3   & \sci{1.1}{-31} & \sci{3.4}{-31} \\
  MMLU      & CLM, Jev         & 0   & 128 & \sci{5.9}{-39} & \sci{3.5}{-38} \\
  MMLU      & CLM, Laya        & 36  & 46  & 0.320          & 0.641 \\
  MMLU      & Jev, Laya        & 120 & 2   & \sci{2.8}{-33} & \sci{1.1}{-32} \\
\end{brandtable}

\begin{brandtable}[Exact McNemar tests for routing and BFCL, Holm-adjusted within each suite\label{tab:mcnemar-other}]{l l C C C C}{Suite, {Left, right}, Left only, Right only, $p$, $p_{\text{Holm}}$}
  Routing & CLEF-Flash, CLM  & 55 & 0  & \sci{5.6}{-17} & \sci{2.8}{-16} \\
  Routing & CLEF-Flash, Jev  & 2  & 3  & 1.000          & 1.000 \\
  Routing & CLEF-Flash, Laya & 47 & 3  & \sci{3.7}{-11} & \sci{1.1}{-10} \\
  Routing & CLM, Jev         & 0  & 56 & \sci{2.8}{-17} & \sci{1.7}{-16} \\
  Routing & CLM, Laya        & 14 & 25 & 0.108          & 0.216 \\
  Routing & Jev, Laya        & 47 & 2  & \sci{4.4}{-12} & \sci{1.7}{-11} \\
  BFCL    & CLEF-Flash, CLM  & 68 & 2  & \sci{4.2}{-18} & \sci{2.1}{-17} \\
  BFCL    & CLEF-Flash, Jev  & 0  & 3  & 0.250          & 0.250 \\
  BFCL    & CLEF-Flash, Laya & 11 & 2  & 0.0225         & 0.0449 \\
  BFCL    & CLM, Jev         & 0  & 69 & \sci{3.4}{-21} & \sci{2.0}{-20} \\
  BFCL    & CLM, Laya        & 7  & 64 & \sci{1.3}{-12} & \sci{5.0}{-12} \\
  BFCL    & Jev, Laya        & 12 & 0  & \sci{4.9}{-4}  & \sci{1.5}{-3} \\
\end{brandtable}

\section{MMLU accuracy by subject}\label{app:mmlu}

\Cref{tab:mmlu-subjects} breaks MMLU accuracy down by subject. Each subject has 20 questions, so one
question moves a score by five points.

\begin{brandtable}[MMLU accuracy (\%) by subject, 20 questions each\label{tab:mmlu-subjects}]{L c c c c}{Subject, CLEF-Flash, Jev, Laya, CLM}
  Abstract algebra          & 90  & 95  & 20 & 40 \\
  Anatomy                   & 75  & 85  & 20 & 30 \\
  Business ethics           & 95  & 80  & 45 & 20 \\
  College computer science  & 90  & 100 & 45 & 40 \\
  Conceptual physics        & 95  & 95  & 35 & 35 \\
  High school psychology    & 95  & 95  & 35 & 25 \\
  International law         & 100 & 100 & 35 & 25 \\
  Machine learning          & 95  & 95  & 25 & 10 \\
  Professional medicine     & 95  & 100 & 35 & 25 \\
  World religions           & 95  & 90  & 50 & 45 \\
\end{brandtable}

\clearpage
\section{Reproducing the analysis}\label{app:reproduce}

Every statistic in this report is computed from the saved per-case rows, so the analysis can be
reproduced offline without starting any model server (\cref{code:reproduce}).

\begin{brandcode}[lang=bash, caption={Recomputing the statistics, notebook, figures, and tables from the saved results\label{code:reproduce}}]
python3 -m pip install -r requirements.txt
python3 scripts/analyze_public_benchmarks.py
python3 scripts/build_results_notebook.py
jupyter nbconvert --to notebook --execute --inplace decision_model_benchmarks_results.ipynb
\end{brandcode}

Rerunning inference for one provider requires its server (or, for Jev, an API key) and replaces that
provider's saved rows (\cref{code:rerun}).

\begin{brandcode}[lang=bash, caption={Rerunning all three runners for one provider while its server is up\label{code:rerun}}]
python3 scripts/run_public_benchmarks.py --providers laya --force
python3 scripts/run_model_routing.py --providers laya --force
python3 scripts/run_tool_selection.py --providers laya --force
\end{brandcode}

\brandbibliography{references}
\brandcta
\end{document}
